%! Characteristics of Radical Functions — Unit 6, Lesson 6.0 — Guided Notes (Answer Key)
\documentclass[10pt]{article}
\usepackage{algebra2-article}
\usepackage{algebra2-key}   % pulls in -boxes; do NOT also load -boxes
\newcommand{\TallMath}[1]{$\displaystyle #1\rule[-1.4em]{0pt}{3.2em}$}
\newcommand{\lgrid}[4]{%
  \draw[step=1, linegray!40, very thin] (#1,#3) grid (#2,#4);
  \draw[->, semithick, charcoal] (#1,0)--(#2,0) node[below right, font=\scriptsize]{$x$};
  \draw[->, semithick, charcoal] (0,#3)--(0,#4) node[left, font=\scriptsize]{$y$};}
% Cube root: pgfmath has no cbrt, so plot the two branches separately.
\newcommand{\cbrtpos}[2]{\draw[very thick, forest, domain=#1:#2, samples=70, smooth]
  plot (\x,{pow(\x,1/3)});}
\newcommand{\cbrtneg}[2]{\draw[very thick, forest, domain=#1:#2, samples=70, smooth]
  plot (\x,{-pow(-1*(\x),1/3)});}
% Vocab answer for the key: bold forest term, then the filled definition on a write-line.
% The leading and trailing \par are required: \ansline ends with \dotfill but does NOT end the
% paragraph, so without them each term label gets pulled onto the previous answer's line.
\newcommand{\vocabans}[2]{%
  \par\noindent\textbf{\textcolor{forest}{#1:}}\\[1pt]\ansline{#2}\par}

\begin{document}

\pageheader{Unit 6, Lesson 6.0}{Guided Notes --- Answer Key}
\namedateperiod

\vspace{0.08in}

\begin{objectivebox}
\textbf{By the end of this lesson, I will be able to\dots}
\begin{itemize}\setlength{\itemsep}{2pt}
  \item find a radical function's \emph{domain} by looking at its \emph{index} and \emph{radicand}, and explain why an \emph{even} index restricts the domain but an \emph{odd} index does not;
  \item read the \emph{domain}, \emph{range}, \emph{endpoint}, \emph{intercepts}, \emph{increasing/decreasing} and \emph{positive/negative} intervals, \emph{extrema}, and \emph{end behavior} off a square root or cube root graph;
  \item explain how $y=\sqrt{x}$ and $y=\sqrt[3]{x}$ are the \emph{inverses} of $y=x^2$ and $y=x^3$, and contrast the two radical families with each other.
\end{itemize}
\end{objectivebox}

\vspace{0.05in}

\begin{vocabbox}
\small
Fill in each term as we build it together.
\par\vspace{2pt}   % \par is required: \vocabans's \noindent is a no-op mid-paragraph
\vocabans{Radical function}{a function whose variable sits underneath a radical, $f(x)=\sqrt[n]{\text{expression in }x}$}
\vocabans{Index}{the small number $n$ in the notch of the radical, telling which root; a plain $\sqrt{\phantom{x}}$ has index $2$}
\vocabans{Radicand}{the expression underneath the radical symbol --- the thing whose root is being taken}
\vocabans{Restricted domain (from the radicand)}{with an \emph{even} index the radicand may not be negative, so the domain is the solution of \emph{radicand} $\ge 0$; an odd index has no restriction}
\vocabans{Endpoint}{the single point where a square root graph begins; it is included in the domain, so it is drawn as a solid dot}
\vocabans{Inverse relationship}{two functions that undo each other, swapping inputs with outputs; their graphs are reflections over the line $y=x$ ($\sqrt{x}$ with $x^2$, $\sqrt[3]{x}$ with $x^3$)}
\end{vocabbox}

\begin{hookbox}
Here are the two graphs this whole unit is built on: $y=\sqrt{x}$ and $y=\sqrt[3]{x}$.
\begin{center}
\begin{tikzpicture}[scale=0.36]
  % y = sqrt(x)
  \begin{scope}
    \lgrid{-2}{9}{-2}{4}
    \begin{scope}
      \clip (-2,-2) rectangle (9,4);
      \draw[very thick, forest, domain=0:9, samples=90, smooth] plot (\x,{sqrt(\x)});
    \end{scope}
    \fill[forest] (0,0) circle (5pt);
    \node[charcoal, font=\scriptsize] at (3.5,-3.2) {$y=\sqrt{x}$};
  \end{scope}
  % y = cbrt(x)
  \begin{scope}[xshift=19cm]
    \lgrid{-8}{8}{-3}{3}
    \begin{scope}
      \clip (-8,-3) rectangle (8,3);
      \cbrtpos{0}{8}
      \cbrtneg{-8}{0}
    \end{scope}
    \node[charcoal, font=\scriptsize] at (0,-4.2) {$y=\sqrt[3]{x}$};
  \end{scope}
\end{tikzpicture}
\end{center}
\begin{itemize}\setlength{\itemsep}{1pt}
  \item The cube root graph runs all the way across the plane. The square root graph \textbf{stops dead} at the origin and never appears to the left of it. Which inputs are missing, and why? \ansline{Every negative input. $\sqrt{x}$ asks for a number that squares to $x$, and no real number squares to a negative --- so there is no point to plot for $x<0$.}
  \item Warm-Up item~1 already told you the reason. What was different about $\sqrt{-36}$ and $\sqrt[3]{-8}$? \ansline{$\sqrt{-36}$ has an \emph{even} index and no real value; $\sqrt[3]{-8}=-2$ because an \emph{odd} index keeps the sign of the radicand.}
\end{itemize}
For every family before this one, we asked \emph{what the graph does}. Radical functions make us ask
something new first: \textbf{where is the graph allowed to exist at all?} The answer is written under
the radical.
\end{hookbox}

\begin{notesbox}{1. What Makes a Function Radical --- and Where It Is Defined}
A \textbf{radical function} has its variable underneath a radical symbol:
\[
  f(x) = \sqrt[n]{\;\text{expression in } x\;} .
\]
Two parts have names, and both matter today:
\begin{itemize}[leftmargin=1.5em, itemsep=3pt]
  \item The \textbf{index} $n$ is the small number in the notch of the radical --- it tells you
        \emph{which} root. A plain $\sqrt{\phantom{x}}$ has index \ans{2} (we just don't write it).
  \item The \textbf{radicand} is everything \ans{underneath} the radical symbol.
\end{itemize}

\vspace{3pt}
\textbf{The rule that runs this unit.} From Warm-Up item~1:
\begin{center}
\renewcommand{\arraystretch}{1.35}
\begin{tabularx}{\linewidth}{@{}l X X@{}}
\hline
\textbf{Index} & \textbf{Can the radicand be negative?} & \textbf{Why} \\
\hline
\textbf{Even} ($\sqrt{\phantom{x}}$, $\sqrt[4]{\phantom{x}}$) & \ans{No} & Nothing real, squared, gives a negative. \\
\textbf{Odd} ($\sqrt[3]{\phantom{x}}$, $\sqrt[5]{\phantom{x}}$) & \ans{Yes} & Cubing \emph{keeps} the sign: $(-2)^3=$ \ans{$-8$}. \\
\hline
\end{tabularx}
\end{center}

\vspace{3pt}
\textbf{The one procedure behind every domain question in this unit:} if the index is
\textbf{even}, set the \ans{radicand} $\ge 0$ and solve. If the index is \textbf{odd}, the domain is
\ans{all real numbers}.
\begin{itemize}[leftmargin=1.5em, itemsep=3pt]
  \item $g(x)=\sqrt{x-3}$: \; solve $x-3\ge 0$, so the domain is \ans{$x\ge 3$, or $[3,\infty)$}.
  \item $h(x)=\sqrt{2x+8}$: \; solve $2x+8\ge 0$, so the domain is \ans{$x\ge -4$, or $[-4,\infty)$}.
  \item $k(x)=\sqrt{5-x}$: \; solve $5-x\ge 0$, so the domain is \ans{$x\le 5$, or $(-\infty,5]$}. \\
        \emph{Notice:} this one runs to the \textbf{left}. A negative in front of the $x$ flips which
        way the graph goes.
  \item $m(x)=\sqrt[3]{x-3}$: \; the index is odd, so the domain is \ans{all reals} --- no work needed.
\end{itemize}
This is genuinely new. A polynomial (Unit~4) was defined \emph{everywhere}. A rational function
(Unit~5) lost a few scattered inputs where the denominator was zero. A square root function loses
an entire \textbf{half} of the number line at once.
\end{notesbox}

\begin{notesbox}{2. The Square Root Parent: $y=\sqrt{x}$}
\begin{center}
\begin{minipage}[c]{0.40\linewidth}
\centering
\renewcommand{\arraystretch}{1.25}
\begin{tabular}{|c|c|}
\hline
$x$ & $y=\sqrt{x}$ \\ \hline
$0$ & \ans{0} \\ \hline
$1$ & \ans{1} \\ \hline
$4$ & \ans{2} \\ \hline
$9$ & \ans{3} \\ \hline
\end{tabular}
\end{minipage}
\begin{minipage}[c]{0.55\linewidth}
\centering
\begin{tikzpicture}[scale=0.46]
  \lgrid{-2}{10}{-2}{4}
  \begin{scope}
    \clip (-2,-2) rectangle (10,4);
    \draw[very thick, forest, domain=0:10, samples=100, smooth] plot (\x,{sqrt(\x)});
  \end{scope}
  \fill[forest] (0,0) circle (4pt);
  \fill[charcoal] (1,1) circle (2.4pt);
  \fill[charcoal] (4,2) circle (2.4pt);
  \fill[charcoal] (9,3) circle (2.4pt);
  \node[charcoal, font=\scriptsize, below left] at (0,0) {endpoint};
\end{tikzpicture}
\end{minipage}
\end{center}

\vspace{2pt}
The graph starts at a single point and travels in \emph{one} direction only. That starting point is
called the \textbf{endpoint}, and it is a \emph{solid} dot --- $x=0$ \emph{is} in the domain.
\begin{enumerate}[leftmargin=1.7em, itemsep=3pt]
  \item \textbf{Domain:} \ans{$[0,\infty)$} \quad \textbf{Range:} \ans{$[0,\infty)$}
  \item \textbf{Endpoint:} $($\ans{0}$,\ $\ans{0}$)$ \quad \textbf{Intercepts:} the graph
        crosses both axes at \ans{$(0,0)$}.
  \item \textbf{Increasing / decreasing:} increasing on \ans{$[0,\infty)$}; it is \emph{never}
        \ans{decreasing}.
  \item \textbf{Extrema:} the endpoint is an \textbf{absolute minimum}, with minimum \emph{value}
        $y=$ \ans{0} at $x=$ \ans{0}. There is \textbf{no} absolute
        \ans{maximum}, because the graph keeps climbing forever.
  \item \textbf{Symmetry:} \ans{none} (it is neither even nor odd).
  \item \textbf{End behavior:} as $x\to+\infty$, $y\to$ \ans{$+\infty$} --- but \emph{slowly}. \\
        And there is no left end at all: the graph does not run off to the left, it
        \ans{stops} at the endpoint. \emph{A radical graph can have only one arm.} This is the
        first family of the year for which ``end behavior'' has only \emph{one} answer.
\end{enumerate}
\end{notesbox}

\begin{notesbox}{3. The Cube Root Parent: $y=\sqrt[3]{x}$}
\begin{center}
\begin{minipage}[c]{0.40\linewidth}
\centering
\renewcommand{\arraystretch}{1.25}
\begin{tabular}{|c|c|}
\hline
$x$ & $y=\sqrt[3]{x}$ \\ \hline
$-8$ & \ans{$-2$} \\ \hline
$-1$ & \ans{$-1$} \\ \hline
$0$ & \ans{0} \\ \hline
$1$ & \ans{1} \\ \hline
$8$ & \ans{2} \\ \hline
\end{tabular}
\end{minipage}
\begin{minipage}[c]{0.55\linewidth}
\centering
\begin{tikzpicture}[scale=0.46]
  \lgrid{-9}{9}{-3}{3}
  \begin{scope}
    \clip (-9,-3) rectangle (9,3);
    \cbrtpos{0}{9}
    \cbrtneg{-9}{0}
  \end{scope}
  \fill[charcoal] (-8,-2) circle (2.6pt);
  \fill[charcoal] (-1,-1) circle (2.6pt);
  \fill[charcoal] (0,0) circle (2.6pt);
  \fill[charcoal] (1,1) circle (2.6pt);
  \fill[charcoal] (8,2) circle (2.6pt);
\end{tikzpicture}
\end{minipage}
\end{center}

\vspace{2pt}
Because the index is \textbf{odd}, negatives are welcome under the radical --- and the whole picture
changes.
\begin{enumerate}[leftmargin=1.7em, itemsep=3pt]
  \item \textbf{Domain:} \ans{all reals, $(-\infty,\infty)$} \quad \textbf{Range:}
        \ans{all reals, $(-\infty,\infty)$}
  \item \textbf{Endpoint:} \ans{none} --- there is nothing to stop the graph, so it has none.
  \item \textbf{Increasing / decreasing:} increasing on \ans{$(-\infty,\infty)$}.
  \item \textbf{Extrema:} \ans{none}. The graph falls forever on the left and climbs forever on
        the right, so there is neither an absolute maximum nor an absolute minimum.
  \item \textbf{Symmetry:} the graph is symmetric about the \ans{origin}, so it is an
        \textbf{odd} function --- check it algebraically: $\sqrt[3]{-8}=$ \ans{$-2$} while
        $-\sqrt[3]{8}=$ \ans{$-2$}, so $f(-x)=$ \ans{$-f(x)$}.
  \item \textbf{End behavior:} as $x\to+\infty$, $y\to$ \ans{$+\infty$}; \; as $x\to-\infty$,
        $y\to$ \ans{$-\infty$}. \emph{Two} arms, like an odd-degree polynomial --- but flattening as
        they go, not steepening.
\end{enumerate}
\end{notesbox}

\begin{notesbox}{4. Where These Two Parents Come From --- the Inverse Relationship}
Warm-Up item~3 already showed it: the table for $y=x^2$ and the table for $y=\sqrt{x}$ are the
\emph{same table with the rows swapped}. Squaring and square-rooting \textbf{undo each other} --- they
are \textbf{inverses} --- and so do cubing and cube-rooting. On a graph, undoing looks like a
\textbf{reflection over the line $y=x$}.
\begin{center}
\begin{tikzpicture}[scale=0.40]
  % y = x^2 (right half) and y = sqrt(x)
  \begin{scope}
    \lgrid{-1}{7}{-1}{7}
    \draw[navy, dashed, semithick] (-1,-1)--(7,7);
    \node[navy, font=\scriptsize, above left] at (6.6,6.6) {$y=x$};
    \begin{scope}
      \clip (-1,-1) rectangle (7,7);
      \draw[very thick, goldacc, domain=0:2.65, samples=80, smooth] plot (\x,{(\x)*(\x)});
      \draw[very thick, forest, domain=0:7, samples=90, smooth] plot (\x,{sqrt(\x)});
    \end{scope}
    \fill[charcoal] (0,0) circle (3pt);
    \node[goldacc, font=\scriptsize, right] at (3.4,6.4) {$y=x^2,\ x\ge 0$};
    \node[forest, font=\scriptsize, above] at (5.6,2.4) {$y=\sqrt{x}$};
  \end{scope}
  % y = x^3 and y = cbrt(x)
  \begin{scope}[xshift=13cm]
    \lgrid{-3}{3}{-3}{3}
    \draw[navy, dashed, semithick] (-3,-3)--(3,3);
    \begin{scope}
      \clip (-3,-3) rectangle (3,3);
      \draw[very thick, goldacc, domain=-1.45:1.45, samples=80, smooth] plot (\x,{(\x)*(\x)*(\x)});
      \cbrtpos{0}{3}
      \cbrtneg{-3}{0}
    \end{scope}
    \node[goldacc, font=\scriptsize, right] at (1.7,2.5) {$y=x^3$};
    \node[forest, font=\scriptsize, below] at (2.6,0.75) {$y=\sqrt[3]{x}$};
  \end{scope}
\end{tikzpicture}
\end{center}
\begin{itemize}[leftmargin=1.5em, itemsep=3pt]
  \item Reflecting over $y=x$ \textbf{swaps the coordinates} of every point: $(4,2)$ on $y=\sqrt{x}$
        came from $($\ans{2}$,\ $\ans{4}$)$ on $y=x^2$. So the two functions also swap their
        \ans{domains} and their \ans{ranges}.
  \item That is the \emph{deeper} reason for today's domain rule. The domain of $y=\sqrt{x}$ is
        $[0,\infty)$ because that is the \ans{range} of $y=x^2$ --- a parabola never produces a
        negative output, so a square root never accepts a negative input.
  \item \textbf{Why only half the parabola?} A full parabola sends \emph{two} inputs to the same
        output ($(-3)^2=3^2=$ \ans{9}), so reflecting all of it would fail the vertical line
        test. To get a function back, we keep only $x\ge 0$. We say $y=x^2$ needs a
        \textbf{restricted domain} to have an inverse.
  \item \textbf{Why does the cube root need no restriction?} Because $y=x^3$ never sends two
        different inputs to the same output: $(-2)^3=$ \ans{$-8$} but $2^3=$ \ans{8}. All of
        it can be reflected, which is exactly why $y=\sqrt[3]{x}$ has domain \ans{all reals}.
\end{itemize}
\emph{Hold on to this.} It is the reason radical functions look the way they do, and Lessons 6.6 and
6.7 will make it a full method.
\end{notesbox}

\begin{notesbox}{5. The Full Feature List --- and How the Two Families Differ}
Now read \emph{every} characteristic off the anchor $f(x)=\sqrt{x+4}-1$.
\begin{center}
\begin{tikzpicture}[scale=0.52]
  \lgrid{-6}{7}{-3}{4}
  \begin{scope}
    \clip (-6,-3) rectangle (7,4);
    \draw[very thick, forest, domain=-4:7, samples=100, smooth] plot (\x,{sqrt((\x)+4)-1});
  \end{scope}
  \fill[forest] (-4,-1) circle (4pt) node[left, font=\scriptsize]{$(-4,-1)$};
  \fill[charcoal] (-3,0) circle (2.6pt) node[above left, font=\scriptsize]{$(-3,0)$};
  \fill[charcoal] (0,1) circle (2.6pt) node[above left, font=\scriptsize]{$(0,1)$};
  \fill[charcoal] (5,2) circle (2.6pt) node[below right, font=\scriptsize]{$(5,2)$};
\end{tikzpicture}
\end{center}
\begin{enumerate}[leftmargin=1.7em, itemsep=3pt]
  \item \textbf{Domain} (set the radicand $x+4\ge 0$): \ans{$[-4,\infty)$} \quad \textbf{Range:}
        \ans{$[-1,\infty)$}
  \item \textbf{Endpoint:} \ans{$(-4,-1)$} --- and notice the domain and range both \emph{start} at
        its coordinates.
  \item \textbf{$x$-intercept / zero:} \ans{$(-3,0)$} \quad \textbf{$y$-intercept:} \ans{$(0,1)$}
  \item \textbf{Increasing / decreasing:} increasing on \ans{$[-4,\infty)$} --- the entire domain.
  \item \textbf{Positive / negative:} $f(x)<0$ on \ans{$[-4,-3)$}; \quad $f(x)>0$ on
        \ans{$(-3,\infty)$}.
        \emph{Careful:} neither interval may include inputs outside the domain.
  \item \textbf{Extrema:} absolute \ans{minimum} of \ans{$-1$} at $x=$ \ans{$-4$}; no
        absolute \ans{maximum}.
  \item \textbf{End behavior:} as $x\to+\infty$, $f(x)\to$ \ans{$+\infty$}. \; As $x\to-\infty$:
        \ans{nothing --- the graph stops at $x=-4$}.
\end{enumerate}

\vspace{4pt}
\textbf{Compare and contrast} the two radical families (this is what makes them one unit but two
pictures):
\begin{center}
\renewcommand{\arraystretch}{1.3}
\begin{tabularx}{\linewidth}{@{}l X X@{}}
\hline
 & \textbf{Square root} $y=\sqrt{x}$ & \textbf{Cube root} $y=\sqrt[3]{x}$ \\
\hline
Index & even & \ans{odd} \\
Domain & \ans{$[0,\infty)$} & \ans{all reals} \\
Range & \ans{$[0,\infty)$} & \ans{all reals} \\
Endpoint & yes, at \ans{$(0,0)$} & \ans{none} \\
Symmetry & none & \ans{origin (odd)} \\
Extrema & absolute \ans{min $0$} & \ans{none} \\
Arms & \ans{one} & \ans{two} \\
Inverse of & $y=x^2$ (restricted) & \ans{$y=x^3$} \\
\hline
\end{tabularx}
\end{center}
And against the families you already know: a \textbf{polynomial} was defined everywhere and had no
edges; a \textbf{rational} function lost scattered inputs and gained asymptotes. A \textbf{radical}
function is the first one whose graph can simply \textbf{begin somewhere and stop}.
\end{notesbox}

\begin{practicebox}
The graph below is $p(x) = \sqrt{4-x}$. Read each characteristic and write it on the line.
\begin{center}
\begin{tikzpicture}[scale=0.5]
  \lgrid{-6}{7}{-2}{4}
  \begin{scope}
    \clip (-6,-2) rectangle (7,4);
    \draw[very thick, forest, domain=-6:4, samples=100, smooth] plot (\x,{sqrt(4-(\x))});
  \end{scope}
  \fill[forest] (4,0) circle (4pt) node[below right, font=\scriptsize]{$(4,0)$};
  \fill[charcoal] (0,2) circle (2.6pt) node[above right, font=\scriptsize]{$(0,2)$};
  \fill[charcoal] (3,1) circle (2.6pt);
\end{tikzpicture}
\end{center}
\begin{enumerate}[leftmargin=1.7em, itemsep=3pt]
  \item Set the radicand $\ge 0$: $4-x\ge 0$, so the \textbf{domain} is \ans{$(-\infty,4]$}. \quad
        \textbf{Range:} \ans{$[0,\infty)$}
  \item \textbf{Endpoint:} \ans{$(4,0)$}. \quad Is $p$ increasing or decreasing on its domain?
        \ans{decreasing}
  \item \textbf{$x$-intercept:} \ans{$(4,0)$} \quad \textbf{$y$-intercept:} \ans{$(0,2)$}
  \item This graph opens to the \textbf{left} while the parent $y=\sqrt{x}$ opens to the right.
        Which part of the equation caused that? \ans{the $-x$ in the radicand}
  \item Does $p$ have an absolute maximum, an absolute minimum, both, or neither?
        \ans{an absolute minimum only ($0$, at $x=4$)}
\end{enumerate}
\end{practicebox}

\begin{teachernote}
\textbf{Pacing.} Sections~1--3 are the core; Section~4 is the conceptual payoff and Section~5 the
assessed skill. If time is short, compress Section~4 to the reflection picture and the domain/range
swap --- but do \emph{not} cut it. It is the ``\,inverse relationship of families\,'' row this
Lesson~0 introduces, and Lessons~6.6--6.7 are built directly on it.

\textbf{The three errors to expect.} \emph{(1)} Writing the domain of $\sqrt{5-x}$ as $x\ge 5$ ---
students solve $5-x\ge0$ without flipping. Warm-Up item~2(c) exists to preempt this; send them back
to it. \emph{(2)} Claiming a cube root graph has a restricted domain ``because it's a radical.'' The
index, not the radical symbol, is what restricts. \emph{(3)} Reporting an end behavior for
$x\to-\infty$ on a square-root graph. Insist on the honest answer: \emph{the graph stops}. This is
the first family of the year where one of the two end-behavior questions has no answer, and students
who do not confront it now will invent a left arm on the Unit~6 test.

\textbf{One more distinction worth saying out loud.} The endpoint is an absolute \emph{minimum} for
$y=\sqrt{x}$ but an absolute \emph{maximum} for a reflected graph like $y=-\sqrt{x}+3$ (homework
item~2). The endpoint is always an extremum on a square-root graph --- which one depends on which way
the arm travels. Cube root graphs, having no endpoint, have no extrema at all.
\end{teachernote}

\end{document}
